\section{Introduction}\label{sec:intro}

A manager who must keep a system above a floor while observing it only partly faces a problem with several independent degrees of freedom that are usually fixed tacitly. Which function of the state is observed? May the same instrument be applied twice in a row, and if not, who remembers what was applied last? If two agencies each control one instrument, how is a disagreement resolved? How often may the regime be reviewed, and how coarse may the monitoring be before the floor is lost? Each of these choices changes the viability kernel, and the changes are not small: on the worked system of this article the number of viable beliefs varies between $111$ and $255$ across four natural information classes, and between $30$ and $255$ once decentralisation is included.

The obstruction calculus developed in \cite{Abaee2026calc} states these problems as common-action problems on beliefs: a belief---a set of states the controller cannot tell apart---is viable when some action keeps every state in it safe and leads to viable beliefs again. The calculus is general, and its general statements are easy to prove; what is hard is to know which of the apparent phenomena are real. Is the ordering between information classes a theorem or an artefact of an example? Does a no-repeat protocol lose viability because of the protocol or because of how the record of the last instrument is kept? Is the loss under decentralisation a property of the voting rule or of the particular belief? Is one-step compatibility of a monitoring partition enough for dynamic viability? These questions are decided by \emph{audits}: exact computations on worked systems in which every claimed inclusion is recomputed from the primitives, every strictness is witnessed, and every closed form is confronted with an independent enumeration. This article is the audit layer of the calculus, rewritten as a research article: the computations are kept, but they are organised around the theorems they establish or refute.

\paragraph{The worked system.}
Two stocks $z_1,z_2\in\{0,1,2,3\}$ must each stay at or above $1$. Instrument $1$ raises $z_1$ by one (capped at $3$) while $z_2$ declines by one; instrument $2$ does the reverse. The corner $(1,1)$ is lost immediately whatever is done, and the other eight safe states are individually viable. Two observation maps are compared---the full state and the aggregate $z_1+z_2$---and two protocols: free choice, and the no-repeat rule that forbids applying the instrument used last period. The four combinations are the \emph{full}, \emph{codex} (full observation, no repeat), \emph{aggregate} and \emph{institutional} (aggregate observation, no repeat) classes. A decentralised variant gives each stock to an agency that votes by a law depending only on its own stock, with the idle instrument applied on disagreement. The same grid supports a hidden-regime variant, a review-timing grid, and a monitoring-design problem, and a continuous companion fibre with two capped instruments gives the calculus its linear-programming face.

\paragraph{Contributions.}
\begin{enumerate}[leftmargin=2em,itemsep=1pt]
\item \emph{A fibrewise belief game and a monotonicity theorem} (Section~\ref{sec:frame}--\ref{sec:mono}). Viability under an observation map, a protocol and a register convention is the greatest fixed point of one operator on beliefs. Kernels are down-closed, stabilise in finitely many steps (three on the worked system), and are monotone along the observation, protocol and horizon axes (Theorem~\ref{thm:mono}); the worked system shows every axis strict and the two middle classes incomparable.
\item \emph{A complete kernel census with closed forms} (Section~\ref{sec:census}). The kernels of the four classes on all $511$ beliefs are $111$, $143$, $119$ and $255$. With two instruments the no-repeat protocol is strict alternation (Lemma~\ref{lem:alt}); the shared-ledger kernel depends on the observation only through the first step and the private-register kernel is a phase-coherence condition (Theorem~\ref{thm:closed}); the two conventions are incomparable, the ledger institutional kernel equals the blind kernel, and the aggregate kernel is realised by the $32$ alternation policies although no memoryless policy sustains any state (Proposition~\ref{prop:stationary}).
\item \emph{Decentralised unanimity factors statewise} (Section~\ref{sec:dec}). A belief is sustained if and only if its states share a sustaining law pair (Theorem~\ref{thm:factor}); $12$ of the $28$ coordinated-viable pairs and $30$ of the $255$ beliefs survive; fixing one agency's vote at the floor moves the kernel between $12$ and $4$ pairs.
\item \emph{Monitoring design as graph structure} (Section~\ref{sec:monitor}). One-step adequacy of a partition is adequacy against an incompatibility graph, with an inclusion--exclusion count and the finite Helly numbers that delimit the graph criterion (Theorem~\ref{thm:adequacy}); the whole set of individually viable states is dynamically viable under a partition if and only if the partition is one-step adequate on it (Theorem~\ref{thm:dynamic}). The census over all $4{,}140$ partitions of the eight viable states gives $1{,}694$ adequate designs with a closed-form count, $16$ coarsest designs, $74$ distinct kernels, and none adequate under the shared ledger.
\item \emph{Regimes, review timing and the continuous fibre} (Sections~\ref{sec:timing}--\ref{sec:continuous}). Mode-blind operation on the explicit hidden-regime model loses exactly the corner; review timing is a max-plus hitting time, with $33$ and $43$ viable cells on a $93$-cell grid and a demonstration that the identity must be stated with the real value and a non-strict inequality; the continuous fibre has crossover $27/5$, margin law $(Y-27/5)/10$, a saddle of value $-1/10$ with a nonconvex gap, and a Helly-tight triple, and the margins are robustness radii.
\item \emph{Verification} (Section~\ref{sec:verification}). An exact Python core recomputes every number from the primitives; the ledger and private kernels are computed by two independent implementations (the generic game engine and the closed forms) that agree on all $511$ beliefs; the verifier runs \NCHECKS{} named checks; the monotonicity theorem is formalised in Lean~4.
\end{enumerate}

\paragraph{Organisation.}
Section~\ref{sec:frame} fixes the worked system, the belief game and the register conventions. Section~\ref{sec:mono} proves the monotonicity theorem. Section~\ref{sec:census} gives the kernel census and the closed forms. Section~\ref{sec:dec} treats decentralised unanimity, Section~\ref{sec:monitor} monitoring design, Section~\ref{sec:timing} regimes and review timing, Section~\ref{sec:continuous} the continuous fibre and the margins. Section~\ref{sec:rules} collects the design rules the audit supports, Section~\ref{sec:verification} describes the verification, Section~\ref{sec:related} places the work, and Section~\ref{sec:conclusion} concludes. The supplementary material contains the master table of all $36$ pairs, the full census tables, the law tables, the timing grids, the benchmark tables, the scaling of the census to larger caps, and the list of checks.
